%
%  tlm_appendix_catalogo_dei_pattern.tex
%
%  Copyright (c) 2026 Roberto Corradini. All rights reserved.
%
%  This file is part of the reversi program
%  http://github.com/rcrr/reversi
%
%  This program is free software; you can redistribute it and/or modify it
%  under the terms of the GNU General Public License as published by the
%  Free Software Foundation; either version 3, or (at your option) any
%  later version.
%
%  This program is distributed in the hope that it will be useful,
%  but WITHOUT ANY WARRANTY; without even the implied warranty of
%  MERCHANTABILITY or FITNESS FOR A PARTICULAR PURPOSE. See the
%  GNU General Public License for more details.
%
%  You should have received a copy of the GNU General Public License
%  along with this program; if not, write to the Free Software
%  Foundation, Inc., 59 Temple Place - Suite 330, Boston, MA  02111-1307, USA
%  or visit the site <http://www.gnu.org/licenses/>.
%

\captionsetup[subfigure]{justification=centering} % Forza tutte le subcaption a essere centrate

\newrobustcmd{\hexmask}[1]{%
    \\[0.5ex] % Va a capo
    \makebox[\textwidth][c]{% Centra perfettamente
        \fontsize{6}{7.2}\selectfont\texttt{0x#1}%
    }
}

% Spostato nella intro.
%\newcounter{myhexcount}

\def\extracthexvalues#1{%
  \setcounter{myhexcount}{1}%
  \foreach \hex in {#1} {%
    \expandafter\xdef\csname hex\Alph{myhexcount}\endcsname{\hex}%
    \stepcounter{myhexcount}%
  }%
}

\ExplSyntaxOn

\NewDocumentCommand{\formatPositions}{m}
{
  \user_regex_format:n { #1 }
}

% Variabili di supporto
\seq_new:N \l_pos_matches_seq
\int_new:N \l_pos_counter_int

\cs_new_protected:Nn \user_regex_format:n
{
  % 1. Resetta il contatore a zero
  \int_zero:N \l_pos_counter_int

  % 2. Trova tutti i match del pattern [a-h][1-8]
  % e salvali nella sequenza \l_pos_matches_seq
  \regex_extract_all:nnN { [a-h][1-8] } { #1 } \l_pos_matches_seq

  % 3. Cicla su ogni elemento trovato
  \seq_map_inline:Nn \l_pos_matches_seq
  {
    % Stampa il comando desiderato con il valore attuale del contatore
    \posannotation{##1}{\int_use:N \l_pos_counter_int}
    % Incrementa il contatore per il prossimo elemento
    \int_incr:N \l_pos_counter_int
    % Aggiungi un piccolo spazio tra i comandi (opzionale)
    \space
  }
}

\ExplSyntaxOff

\newcommand{\drawSubfigure}[4]{% 
  \begin{subfigure}[b]{0.23\textwidth}
    \centering
    \scalebox{0.45}{
      \begin{othelloboard}{1}
        \expandafter\formatPositions\expandafter{#4}%
      \end{othelloboard}
    }
    \caption{$T_{#1}$ : #2 \hexmask{#3}}
  \end{subfigure}
}


\section{Catalogo dei Pattern}

In questa appendice vengono visualizzati i pattern principali utilizzati dal sistema. Per ogni pattern sono mostrate le 8 trasformazioni del gruppo $D_4$. I numeri all'interno delle celle indicano l'indice di packing progressivo (0, 1, 2, ...) all'interno del pattern, permettendo di tracciare visivamente la permutazione delle celle sotto ogni trasformazione.

I Pattern sono categorizzati secondo i Sottogruppi di $D_4$ e propongono un set esaustivo delle casistiche esistenti.

\begin{enumerate}
  \item \textbf{Livello 8 (Cima):} Rappresenta il gruppo $D_4$ stesso. Contiene tutte le 8 simmetrie.
  \item \textbf{Livello 4 (Intermedio):} Contiene i sottogruppi massimali. Notiamo $C_4$ (rotazioni pure) e i due gruppi di Klein $V_4$. Ogni gruppo di questo livello contiene esattamente 3 sottogruppi di ordine 2.
  \item \textbf{Livello 2 (Atomi):} Composti dall'identità e da un elemento invertibile (una riflessione o la rotazione di $180^\circ$). Sono i "mattoni" fondamentali del gruppo.
  \item \textbf{Livello 1 (Base):} Il sottogruppo banale $\{e\}$.
\end{enumerate}

\begin{table}[h]
  \centering
  \newcolumntype{M}{>{\ttfamily}l}
  \begin{tabular}{lclMlll}
    \toprule
    \textbf{Tipo} & \textbf{Livello} & \textbf{Pattern} & \multicolumn{1}{l}{\textbf{Fingerprint}} & \textbf{Simmetrie} & \textbf{Orbite} & \textbf{Stab} \\
    \midrule
    0  & 1 & ELLE   & [I0, T1, T2, T3, T4, T5, T6, T7] & ${\{e\}}$         & 8 & 1 \\
    1  & 2 & SNAKE  & [I0, T1, S0, S1, T4, T5, S4, S5] & ${\{e, R_{180}\}}$ & 4 & 2 \\
    2  & 2 & EDGE   & [I0, T1, T2, T3, S0, S1, S2, S3] & ${\{e, S_v\}}$     & 4 & 2 \\
    3  & 2 & DIAG3  & [I0, T1, T2, T3, S1, S2, S3, S0] & ${\{e, S_{d2}\}}$  & 4 & 2 \\
    3D & 2 & DOTA1  & [I0, T1, T2, T3, I1, I2, I3, I0] & ${\{e\}}$          & 4 & 2 \\
    4  & 2 & TAU    & [I0, T1, T2, T3, S2, S3, S0, S1] & ${\{e, S_h\}}$     & 4 & 2 \\
    5  & 2 & MACE   & [I0, T1, T2, T3, S3, S0, S1, S2] & ${\{e, S_{d1}\}}$  & 4 & 2 \\
    6  & 4 & COREA  & [I0, S0, S0, S0, T4, S4, S4, S4] & ${C_4}$            & 2 & 4 \\
    7  & 4 & BARBEL & [I0, T1, S0, S1, S1, S0, S1, S0] & ${V_{4(d)}}$        & 2 & 4 \\
    7D & 4 & DIAG8  & [I0, T1, S0, S1, I3, I0, I1, I2] & ${\{e, R_{180}\}}$  & 2 & 4 \\
    8  & 4 & RCT2X4 & [I0, T1, S0, S1, S0, S1, S0, S1] & ${V_{4(m)}}$        & 2 & 4 \\
    9  & 8 & CORE   & [I0, S0, S0, S0, S0, S0, S0, S0] & ${D_4}$            & 1 & 8 \\
    \bottomrule
  \end{tabular}
  \caption{Categorizzazione dei Pattern}
\end{table}

\newpage

% Inizia il ciclo per ogni riga del CSV mdp.csv
\DTLforeach{master_data_patterns}{
  \mdpName=name,
  \mdpMask=mask,
  \mdpNSquares=n_squares,
  \mdpNConfigurations=n_configurations,
  \mdpNInstances=n_instances,
  \mdpNStabilizers=n_stabilizers,
  \mdpCells=cells,
  \mdpTMasks=tmasks,
  \mdpMaskIndexes=mask_indexes,
  \mdpUniqueMasks=unique_masks,
  \mdpUniqueMaskIndexes=unique_mask_indexes,
  \mdpTr=tr,
  \mdpAt=at,
  \mdpSf=sf,
  \mdpCellsTA=cells_t0,
  \mdpCellsTB=cells_t1,
  \mdpCellsTC=cells_t2,
  \mdpCellsTD=cells_t3,
  \mdpCellsTE=cells_t4,
  \mdpCellsTF=cells_t5,
  \mdpCellsTG=cells_t6,
  \mdpCellsTH=cells_t7,
  \mdpFingerprint=fingerprint,
  \mdpType=type}{%

  \StrSubstitute{\mdpCells}{:}{, }[\mdpCellsCleaned]
  \StrSubstitute{\mdpTMasks}{:}{, }[\mdpTMasksCleaned]
  \StrSubstitute{\mdpMaskIndexes}{:}{, }[\mdpMaskIndexesCleaned]
  \StrSubstitute{\mdpUniqueMasks}{:}{, }[\mdpUniqueMasksCleaned]
  \StrSubstitute{\mdpUniqueMaskIndexes}{:}{, }[\mdpUniqueMaskIndexesCleaned]
  \StrSubstitute{\mdpTr}{:}{, }[\mdpTrCleaned]
  \StrSubstitute{\mdpAt}{:}{, }[\mdpAtCleaned]
  \StrSubstitute{\mdpSf}{:}{, }[\mdpSfCleaned]
  \StrSubstitute{\mdpFingerprint}{:}{, }[\mdpFingerprintCleaned]
  
  \StrSubstitute{\mdpTMasks}{:}{,}[\mdpTMasksComma]
  \expandafter\extracthexvalues\expandafter{\mdpTMasksComma}
  
  \subsection{Pattern \mdpName}
  
  \begin{tcolorbox}[sharp corners, fontupper=\ttfamily\small, colback=white]
    \noindent Pattern: name = \mdpName, mask = 0x\mdpMask\\
    {[} n\_squares = \mdpNSquares, n\_configurations = \mdpNConfigurations,
    n\_instances = \mdpNInstances, n\_stabilizers = \mdpNStabilizers, type = \mdpType { ]}\\[1.0ex]
    \begin{tabularx}{\textwidth}{@{} p{3.5cm} X @{}}
      Cells:               & {[}\mdpCellsCleaned{]}\\[0.5ex]
      Transformed masks:   & \mdpTMasksCleaned\\[0.5ex]
      Mask indexes:        & {[}\mdpMaskIndexesCleaned{]}\\[0.5ex]
      Unique masks:        & \mdpUniqueMasksCleaned\\[0.5ex]
      Unique mask indexes: & {[}\mdpUniqueMaskIndexesCleaned{]}\\[0.5ex]
      Transf. functions:   & {[}\mdpTrCleaned{]}\\[0.5ex]
      Anti-transf. f.:     & {[}\mdpAtCleaned{]}\\[0.5ex]
      Symmetry functions:  & {[}\mdpSfCleaned{]}\\[0.5ex]
    \end{tabularx}
  \end{tcolorbox}

  \begin{figure}[htbp]
    \centering
    \drawSubfigure{0}{ro000}{\hexA}{\mdpCellsTA} \hfill
    \drawSubfigure{1}{ro090}{\hexB}{\mdpCellsTB} \hfill
    \drawSubfigure{2}{ro180}{\hexC}{\mdpCellsTC} \hfill
    \drawSubfigure{3}{ro270}{\hexD}{\mdpCellsTD}
    \vspace{0.3cm}
    \drawSubfigure{4}{fvert}{\hexE}{\mdpCellsTE} \hfill
    \drawSubfigure{5}{fh1a8}{\hexF}{\mdpCellsTF} \hfill
    \drawSubfigure{6}{fhori}{\hexG}{\mdpCellsTG} \hfill
    \drawSubfigure{7}{fa1h8}{\hexH}{\mdpCellsTH}
    \caption{Istanze e indici di permutazione del pattern \mdpName.}
  \end{figure}
  
  \newpage
}
